% !TEX root = ../main.tex
% called by main.tex
\chapter{Introducción al software CAD}
\label{ch:cad-intro}

\section {Autodesk Fusion}
Autodesk Fusion es un programa de diseño CAD integrado en la nube desarrollado por Autodesk. Este programa es donde se han desarrollado a día de hoy la mayoría de los diseños para el club de vuelo, ya que permite una integración rápida y ofrece una barrera de entrada a nuevos usuarios menor que otros softwares comerciales más extendidos en la industria (como CATIA V5). 
\par Las características por las que hemos elegido Fusion frente a otros son:
\begin{itemize}
    \item Es gratis, gracias a la cuenta de \texttt{alumnos.upm}
    \item Permite colaboración en la nube, evitando problemas de tranferencias de archivos
    \item Ofrece numerosos servicios, desde diseño paramétrico de componentes, hasta diseño de circuitos impresos o renderizado y animación.
    \item Por defecto cuenta con numerosas ayudas para nuevos usuarios que facilitan su introducción al CAD.
\end{itemize}

\subsection{Instalación de Fusion}
Esta parte está todavía en trabajo, me gustaría añadir una pequeña guía e imágenes del proceso. \todo{Instalación de Fusion}

\section{CATIA V5}
\newcommand{\catia}{CATIA V5 }
\catia es otro programa de diseño CAD mucho más arraigado en los sectores aeronáuticos e industrial, desarrollado por Dassault Systémes. En cuanto a las capacidades que ofrece, es un software muy potente, con una documentación muy extensa que permite aprender las funcionalidades más complejas a usuarios con el tiempo y la paciencia para buscarla. 
\par Un pequeño inconveniente de este software es que tiene tanta herencia (CATIA se desarrolló por primera vez en 1982) que no se ha decidido actualizar elementos fundamentales para facilitar la introducción de nuevos usuarios al CAD, como la interfaz de usuario, la funcionalidad o la distribución de los distintos entornos de trabajo. Siendo el objetivo de esta guía facilitar el diseño con la impresión 3D como fin, es importante recalcar que no recomiendo utilizar \catia como la entrada al diseño CAD, pero sí que será perfectamente adecuado para trabajar. 

\section{Prusa Slicer}
\todo{sección de PrusaSlicer}

\section{Ultimaker Cura}
\todo{sección de Cura}